\answer{ex:03:1}{$35\to17.5\mV$ com $g_E=g_L$; $56\to40\mV$ com $g_E=4g_L$; a subtração daria $38.5\mV$: a derivação divide $g_E$ por $3$. A janela se estreita pela metade ($\tau_{\text{ef}}$ cai $2$ vezes).}
\answer{ex:03:2}{$\operatorname{tr}=\frac{w_{EE}-1}{\tau_E}-\frac1{\tau_I}$, $\det=\frac{1-w_{EE}+w_{EI}w_{IE}}{\tau_E\tau_I}$; na fronteira $\omega=\sqrt{\det}$, $\tau_I=20\ms$, período de $73\ms$.}
\answer{ex:03:3}{$d_c=\tau_E\arccos(-1/G)/\sqrt{G^2-1}$, período $2\pi\tau_E/\sqrt{G^2-1}\in(2d_c,4d_c)$. $G=2$: $d_c=12.1\ms$, período de $36\ms$; $G=5$: $d_c=3.6\ms$, período de $12.8\ms$: exatamente a faixa dos ritmos rápidos.}
\answer{ex:03:4}{(a)~Em $I_2=\beta I_1=2$ subindo e $I_2=I_1/\beta=0.5$ descendo, um laço de largura $1.5$. (b)~Em $\tau\ln10/(\beta-1)=2.3\tau$ a cada ordem decimal de $\varepsilon$.}
\answer{ex:03:5}{Três intermediários, $d_k=5$, $10$, $15\ms$: os vetos devem cobrir $15\ms$, $5\ms$ cada um.}
\answer{ex:03:6}{$n+M+1=3+4+1=8$ neurônios. Dois: \nt{GABA} $=[x+y+z\ge2]$, saída $=[x+y+z-2\,\nt{GABA}\ge1]$.}
\answer{ex:03:7}{$\binom84=70<100\le\binom94=126$: $9$ intermediários (teorema de Sperner). Sem conexões diretas, $100$: o posto de $\beta(J-I)$ é $n$.}
\answer{ex:03:8}{$h_c=\frac1{2w_{EE}}+\frac{w_{EI}w_{IE}-w_{EE}}{4w_{EE}^2}=\frac7{18}\approx0.39$; a simulação dá a troca de sinal entre $h=0.38$ e $0.40$.}
